\slajdid{10-s040}
\begin{frame}[label=10-s040]
\gusttekst\setlength{\abovedisplayskip}{3pt}\setlength{\belowdisplayskip}{3pt}\setlength{\abovedisplayshortskip}{2pt}\setlength{\belowdisplayshortskip}{2pt}\begin{columns}[c,onlytextwidth]\column{.29\textwidth}\centering\begin{tikzpicture}[x=.62cm,y=.8cm,font=\sitneoznake,>=Latex]\ctikzset{bipoles/length=.6cm}\draw(-1.8,1)node[ocirc]{}node[left]{$V_{I1}$}--(-1,1);\begin{scope}[shift={(0,1)}]\draw(-1,0)--(-.3,0);\draw(-.3,-.3)--(-.3,.3)--(.3,0)--cycle;\draw(.3,0)--(1,0);\draw(.3,-.3)--(.3,.3);\end{scope}\draw(1,1)--(1.4,1);\node[below]at(0,0.85){D1};\draw[<->](-.5,1.55)--(.5,1.55);\draw(-.5,1.4)--(-.5,1.7);\draw(.5,1.4)--(.5,1.7);\node[above]at(0,1.7){$+\ V_{D1}$};\draw[->](-1.7,1.35)--(-1.1,1.35)node[midway,above]{$I_{D1}$};\draw(-1.8,-1)node[ocirc]{}node[left]{$V_{I2}$}--(-1,-1);\begin{scope}[shift={(0,-1)}]\draw(-1,0)--(-.3,0);\draw(-.3,-.3)--(-.3,.3)--(.3,0)--cycle;\draw(.3,0)--(1,0);\draw(.3,-.3)--(.3,.3);\end{scope}\draw(1,-1)--(1.4,-1);\node[below]at(0,-1.15){D2};\draw[<->](-.5,-0.44999999999999996)--(.5,-0.44999999999999996);\draw(-.5,-0.6)--(-.5,-0.30000000000000004);\draw(.5,-0.6)--(.5,-0.30000000000000004);\node[above]at(0,-0.30000000000000004){$+\ V_{D2}$};\draw[->](-1.7,-0.65)--(-1.1,-0.65)node[midway,above]{$I_{D2}$};\draw(1.4,1)--(1.4,-1);\draw(1.4,0)--(2.8,0)node[ocirc]{}node[right]{$V_O$};\draw(2.1,0)to[R,l=$R$](2.1,-1.8)node[ground]{};\draw[->](1.6,-.5)--(1.6,-1.1);\node[left]at(1.6,-1.55){$I_R$};\end{tikzpicture}
\column{.69\textwidth}
Posmatrajmo situaciju kada je na primer $V_{I1}\gg V_{\gamma D}$ dok je $V_{I2}$ i dalje manje od $V_{\gamma D}$. Ako bi sada i dalje pretpostavili da diode ne vode izlazni napon bi bio nula, što znači da je na katodi diode D1 napon jednak nuli. Na anodi diode D1 je ulazni napon pa bi napon na diodi bio $V_{D1}=V_{I1}\gg V_{\gamma D}$. Znači dioda za koju smo smatrali da ne vodi, da je zakočena ima sve uslove da provodi. Pogrešna polazna pretpostavka.
\par\smallskip Smatrajno sada da dioda D1 vodi. U tom slučaju po konturi diode D1 i otpornika R
\[V_{I1}-V_{D1}-V_O=0\]
Kako smatramo da dioda vodi
\[V_{I1}-V_D=V_O\]
odnosno
\[I_R=\frac{V_{I1}-V_D}{R}=I_D>0\]
pošto je pretpostavka $V_{I1}\gg V_{\gamma D}$.
\par\smallskip Struja kroz didodu je pozitivna, što se slaže sa našom polaznom pretpostavkom da dioda vodi. Napon na izlazu je $V_O=V_{I1}-V_D$. Dioda D2 ne vodi, inverzno je polarizovana pošto je $V_{D2}=V_{I2}-V_O=0-(V_{I1}-V_D)<0$.
\end{columns}
\end{frame}
